\documentclass[main.tex]{subfiles}

\begin{document}

\section{Results}
\label{sec:results}

\subsection{Internal Layer Behavior}

Before evaluating the complete method, we examined whether internal CNN
layers provided useful candidate metric views under the validation protocol.
Each layer was evaluated independently with distance-weighted kNN, and
validation macro-F1 was used only to characterize and rank the candidate
views. Because the backbones expose different numbers of hooks, layer
positions were normalized within each architecture and grouped as early,
middle, deep, or final.

Figure~\ref{fig:internal_layers}A reports standalone validation performance
after first averaging layers within each depth group and context. Mean
macro-F1 increased from 0.8920 for early views to 0.9666 for deep views,
while final views reached 0.9582. Deep and final representations were strong
on average, but early and middle views also provided useful validation
signals. Figure~\ref{fig:internal_layers}B shows that these earlier views
also appeared in a subset of the validation-selected prefixes. These
diagnostics support treating internal layers as candidate metric views; they
are not test-set performance claims or clinical explanations.

\begin{figure*}[!t]
  \centering
  \includegraphics[width=\textwidth]{fig02_internal_layer_diagnostics}
  \caption{Internal-layer diagnostics over 45 brain-MRI dataset-backbone
  contexts. Panel A shows standalone validation macro-F1 by normalized CNN
  depth group. Layer scores are first averaged within each depth group and
  context so that architectures contribute equally. Panel B counts the
  selected layer views appearing in validation-selected HMDF-kNN prefixes.
  These diagnostics describe validation behavior and are not test-set
  performance claims.}
  \label{fig:internal_layers}
\end{figure*}

\subsection{Selected HMDF-kNN Configurations}

Figure~\ref{fig:selected_configuration} characterizes the complete
configurations selected by validation. HMDF-kNN selected one layer in 6
contexts, two in 11, three in 15, and four in 13. The neighborhood size was
$k=1$ in 41 contexts. Uniform, deterministic score-based, and finite
reproducible Dirichlet candidate weights were selected in 7, 2, and 36
contexts, respectively. These frequencies show that validation controls the
number of metric views, the neighborhood size, and the distance-weight rule
rather than applying one fixed configuration to every dataset-backbone
context. They describe classifier behavior and should not be interpreted as
clinical explanations of tumor appearance.

\begin{figure*}[!t]
  \centering
  \includegraphics[width=\textwidth]{fig03_selected_configuration}
  \caption{Validation-selected HMDF-kNN configurations over 45 brain-MRI
  contexts. Panel A shows the selected ranked-prefix length, Panel B the
  selected kNN neighborhood size, and Panel C the selected distance-weight
  family. Score-based weights group deterministic validation-score weighting
  rules; Dirichlet groups the finite reproducible candidate set evaluated for
  each prefix.}
  \label{fig:selected_configuration}
\end{figure*}

\subsection{Design Ablation of the HMDF-kNN Aggregation Strategy}

Table~\ref{tab:ablation} studies the internal aggregation design rather than
context-level wins or losses. The best validation-ranked single layer reached
0.9560 mean test macro-F1. Uniform all-layer distance fusion increased the
mean to 0.9582, indicating that multilayer aggregation can be beneficial.
However, blindly aggregating every available layer remained below the
validation-ranked prefix variants and the complete HMDF-kNN configuration.
Ranked-prefix and greedy constructions reached 0.9599--0.9606, while the
complete method reached 0.9616 using an average of 2.78 layer views.

The pattern suggests that the aggregate result is not explained only by
adding more layers. Instead, it is consistent with selecting and aggregating
complementary layer-wise distance spaces under validation control. Because
the variants may jointly change layer selection, aggregation, and $k$, this
table should not be interpreted as an isolated causal test of the weight
family.

\begin{table*}[t]
\caption{Design ablation of HMDF-kNN aggregation strategy.}
\label{tab:ablation}
\centering
\footnotesize
\setlength{\tabcolsep}{5pt}
\begin{tabular}{p{3.5cm}p{4.7cm}p{3.8cm}rr}
\toprule
Variant & Aggregation strategy & Layer selection & Mean layers & Mean test macro-F1 \\
\midrule
Best single layer & None & Best validation-ranked layer & 1.00 & 0.9560 \\
All-layer uniform distance fusion & Uniform distance aggregation & All layers & 5.67 & 0.9582 \\
Ranked-prefix uniform fusion & Uniform distance aggregation & Validation-ranked prefix & 2.49 & 0.9599 \\
Ranked-prefix score-weighted fusion & Score-based distance aggregation & Validation-ranked prefix & 2.49 & 0.9599 \\
Greedy score-weighted fusion & Score-based distance aggregation & Greedy selected subset & 2.56 & 0.9604 \\
Greedy uniform fusion & Uniform distance aggregation & Greedy selected subset & 2.60 & 0.9606 \\
HMDF-kNN & Validation-selected weighted distance aggregation & Validation-ranked prefix & 2.78 & 0.9616 \\
\bottomrule
\end{tabular}
\par\vspace{2pt}
\footnotesize Results are averaged over 45 brain-MRI contexts. Rows summarize complete validation-selected design variants and may jointly differ in layer selection, aggregation, and $k$; they are not an isolated test of the weight family.
\end{table*}


\subsection{Aggregate Comparison by Method Family}

After characterizing the internal behavior and aggregation design of
HMDF-kNN, we compared the complete frozen method against external reference
families. Figure~\ref{fig:block_macro_f1} reports mean test macro-F1 over the
45 brain-MRI contexts. The softmax head reached 0.9326,
validation-selected final-embedding classifiers reached 0.9470, multilayer
reference methods reached 0.9550, and HMDF-kNN reached 0.9616. These values
show a favorable aggregate progression toward representation-aware post-hoc
methods, but do not imply uniform superiority in every context.

\begin{figure*}[!t]
  \centering
  \begin{minipage}{0.64\textwidth}
    \centering
    \includegraphics[width=0.76\linewidth]{fig04_block_macro_f1}
    \caption{Mean test macro-F1 by method family over 45 brain-MRI
    dataset-backbone contexts. Error bars indicate 95\% bootstrap confidence
    intervals across contexts. For each non-proposed family and context, the
    reported method is selected by validation macro-F1. HAM10000 is excluded.}
    \label{fig:block_macro_f1}
  \end{minipage}
\end{figure*}

\subsection{Complete Aggregate Method Comparison}

Table~\ref{tab:complete_methods} expands the family-level summary to the
specific evaluated methods. The strongest non-proposed aggregate results
were obtained by the multiview references MAXVAR-GCCA and MvDA, with mean
macro-F1 values of 0.9515 and 0.9504. Among the individual final-embedding
classifiers, kNN reached 0.9469. HMDF-kNN obtained the highest listed mean
accuracy, macro-F1, and balanced accuracy at 0.9683, 0.9616, and 0.9597,
respectively. These are averages across dataset-backbone contexts and should
be interpreted as aggregate evidence rather than universal context-level
superiority.

\begin{table*}[t]
\caption{Complete aggregate method comparison grouped by method family.}
\label{tab:complete_methods}
\centering
\scriptsize
\setlength{\tabcolsep}{4.2pt}
\begin{tabular}{p{3.1cm}p{7.1cm}rrrr}
\toprule
Method family & Method & $n$ & Acc. & Macro-F1 & Bal. acc. \\
\midrule
Direct CNN baseline & Softmax head & 45 & 0.9386 & 0.9326 & 0.9338 \\
\addlinespace[2pt]
Final-embedding classifiers & Validation-selected final classifier & 45 & 0.9534 & 0.9467 & 0.9499 \\
 & kNN & 45 & 0.9537 & 0.9469 & 0.9516 \\
 & Linear SVM & 45 & 0.9395 & 0.9323 & 0.9330 \\
 & Random forest & 45 & 0.9300 & 0.9219 & 0.9126 \\
 & Logistic regression & 45 & 0.9257 & 0.9185 & 0.9259 \\
 & Diagonal GMM & 45 & 0.9191 & 0.9148 & 0.9102 \\
 & XGBoost & 45 & 0.9199 & 0.9096 & 0.9003 \\
 & Gaussian naive Bayes & 45 & 0.9130 & 0.9055 & 0.8894 \\
\addlinespace[2pt]
Multilayer controls & Raw concatenation + linear & 45 & 0.9409 & 0.9347 & 0.9305 \\
 & Concatenation + PCA + linear & 45 & 0.9412 & 0.9352 & 0.9310 \\
 & Uniform layer soft vote & 45 & 0.9224 & 0.9121 & 0.9005 \\
 & Uniform kernel SVM & 45 & 0.9437 & 0.9369 & 0.9341 \\
\addlinespace[2pt]
Literature multilayer references & MLCFF-style fusion (Fradi et al., 2021)$^\dagger$ \cite{fradi2021multilayer} & 45 & 0.9416 & 0.9345 & 0.9314 \\
 & Head2Toe-style fusion (Evci et al., 2022)$^\dagger$ \cite{evci2022head2toe} & 45 & 0.9352 & 0.9282 & 0.9258 \\
 & EasyMKL$^\dagger$ & 45 & 0.9376 & 0.9307 & 0.9289 \\
\addlinespace[2pt]
Multiview references & MAXVAR-GCCA$^\dagger$ & 45 & 0.9584 & 0.9515 & 0.9496 \\
 & GMLDA$^\dagger$ & 45 & 0.9426 & 0.9360 & 0.9320 \\
 & MvDA$^\dagger$ & 45 & 0.9573 & 0.9504 & 0.9462 \\
\addlinespace[2pt]
Metric/prototype references & NCA + kNN$^\dagger$ & 45 & 0.9486 & 0.9415 & 0.9408 \\
 & KMEx (Gautam et al., 2024) \cite{gautam2024kmex} & 45 & 0.9383 & 0.9328 & 0.9282 \\
\addlinespace[2pt]
This work & HMDF-kNN & 45 & 0.9683 & 0.9616 & 0.9597 \\
\bottomrule
\end{tabular}
\par\vspace{2pt}
\scriptsize Each row averages 45 brain-MRI dataset-backbone contexts. Hyperparameters are selected by validation macro-F1 within each method and context. Individual final-embedding rows select the best hyperparameter setting for that classifier. $^\dagger$Stage-level adaptation over the saved CNN embedding views; complete configurations and per-context outputs are provided in the supplementary artifacts.
\end{table*}


Methods whose principal reference is from 2020 onward include author and
year directly in the method name for clarity. Stage-level adaptation details,
hyperparameter grids, and complete per-context outputs are documented in the
supplementary artifacts. The public code location remains a submission task:
\texttt{[GitHub repository placeholder]}.

\FloatBarrier
\clearpage
\subsection{Context-Level Robustness}

Context-level analyses were used as robustness checks rather than as the
primary narrative. The aggregate advantage of HMDF-kNN was not uniform
across all dataset-backbone contexts; most paired comparisons were favorable
or practically close under the 0.005 macro-F1 margin, while only a small
number favored the validation-selected reference method. Paired bootstrap
differences were statistically resolved in only two contexts after Holm
correction, and none significantly favored a reference method among the 38
contexts with aligned prediction artifacts. This supports a cautious
aggregate interpretation rather than a claim of universal superiority.
Dataset-level summaries, the complete context heatmap, numerical and
statistical outcome tables, class-level confusion matrices, and the
qualitative fused-distance visualization are provided as supplementary
artifacts.

\textit{External-domain control:} HAM10000 was evaluated only to stress-test
the post-hoc procedure outside brain MRI. It is excluded from brain-MRI
averages, context counts, and inferential tests; the main claims remain
restricted to the five brain-MRI datasets.

\end{document}
